\section{Original local canonical principalization with the same smooth Cartesian stages}

Let $X^n$ be an arbitrary complex manifold, $I$ its original coherent
functional ideal, nonzero along every point, and $E_0$ an arbitrary
original ordered SNC boundary. The original datum is
$M_0=(X,X,E_0,I,1)$. The statements retain the actual canonical terminal
towers furnished by the local driver, rather than only their resulting
principal ideals. No compactness of $Z(I)$ is assumed.

\begin{lemma}[Whole geometry of an actual ambient terminal tower]
\label{hir:original-local-canonical-principalization-terminal-geometry}
\lean{Hironaka.HironakaOriginalLocalCanonicalPrincipalization.terminal_principalization_geometry}
\uses{hir:principalization-boundary}
\leanok
Suppose $M=(X,X,E,J,1)$ is an actual invariant-domain datum and $T$ an
actual terminal geometric test tower for $M$. If the whole original zero
locus of $J$ is $F$, then every actual centre lies over $F$, the tower
is admissible for $(X,E,F)$, and its entire ordered boundary $E_T$ is SNC.
Its actual total exceptional divisor $D_T$ satisfies
\[
 T^*J=\mathcal O(-D_T),\qquad
 |D_T|=T^{-1}F=\bigcup(E_T.\operatorname{drop}|E|).
\]
The entire differential exceptional locus lies in $|D_T|$. The same
projection is holomorphic, proper and surjective, and is an isomorphism
over $X\setminus F$.
\end{lemma}
\begin{proof}
Coherence and marking one identify the whole cosupport of the ambient
datum with the whole zero locus of $J$, hence with $F$. Actual testness
confines every centre over that original cosupport and gives admissibility.
The invariant domain is preserved by the genuine marked transforms, so
the entire final ordered boundary is SNC. Actual terminality gives the
whole principalization identity $T^*J=\mathcal O(-D_T)$. The zero locus
of the whole coherent pullback is the actual inverse image of its whole
zero locus, which proves $|D_T|=T^{-1}F$. The actual boundary bookkeeping
identifies this support with precisely the entries born after all the
original $|E|$ positions, with none of those positions deleted. Every
differential exceptional point of an actual finite blowup tower belongs
to this support. The actual blowup sequence is holomorphic, proper and
surjective, and is an isomorphism away from its exceptional image. That
image lies in $F$ by the proved support equality, giving the stated
isomorphism over its complement. This lemma's actual terminal-tower
hypothesis is internally discharged in the following existence statements.
\end{proof}

\begin{theorem}[A prescribed relatively compact open retains its actual canonical tower and every smooth source]
\label{hir:original-local-canonical-principalization-prescribed-open}
\lean{Hironaka.HironakaOriginalLocalCanonicalPrincipalization.exists_canonical_boundary_principalization_of_subset_isCompact}
\uses{hir:original-local-canonical-principalization-terminal-geometry,hir:resolution-relatively-compact,hir:original-canonical-boundary-principalization-smooth}
\leanok
Let $U\subset X$ be a prescribed original open contained in a compact
set $L$. For the unchanged whole restricted datum
$M=M_0|_U$, original geometry internally produces an actual finite
canonical terminal tower $T:U\leftarrow D$ for the actual
$\operatorname{canonicalInvariant}(n)$. It retains the entire geometry
of the preceding lemma, with
$F=U\cap Z(I)$, $J=I|_U$ and the entire restricted original boundary
$E_0|_U$. In particular the support of its total exceptional divisor is
the union of the entries after the original $|E_0|$ positions.

The choice of this same $D,T$ precedes \emph{every} genuine submersion
$b:Y^m\to U$. For each such $b$ it internally produces an actual
finite tower $Q:Y\leftarrow C$ and actual upper submersion $g:C\to D$
in a holomorphic Cartesian square, with equal literal lengths and
$Tg=bQ$. The intrinsic source datum is the unchanged whole marked
pullback $N=b^*M$; it has the entire ordered inverse-image boundary,
the actual whole pulled ideal and marking one. The same $Q$ is canonical
and terminal for $\operatorname{canonicalInvariant}(m)$ on $N$.
Every source centre lies over $b^{-1}F$, and the whole source geometry
of the preceding lemma holds, including
\[
 E_Q=E_T.\operatorname{map}(g^{-1}),\quad
 Q^*(b^*J)=\mathcal O(-D_Q),\quad
 |D_Q|=Q^{-1}(b^{-1}F).
\]
The source born support is the union after the same original $|E_0|$
positions. No source compactness, source resolving history or comparison
certificate is supplied.
\end{theorem}
\begin{proof}
Original coherence, nonzero-alongness and SNC give the actual invariant
domain of $M_0$. Apply the stronger relatively compact canonical driver
to the original $U\subset L$, retaining both its actual canonical tower
and its terminal geometric resolution. Its restricted datum is exactly
$M$, not a replacement ideal or boundary. Open pullback preserves the
actual domain, and the whole zero-locus pullback formula gives
$J.\operatorname{zeroLocus}=F$. Apply the preceding lemma to this same
chosen $T$. The original boundary restriction is a literal list map,
so it has exactly $|E_0|$ positions.

Now fix an arbitrary actual submersion $b$ after this target choice.
Genuine stepwise Cartesian blowup construction produces $Q,g$ from the
same $T$, retaining every actual empty centre. Its actual universal
property gives the holomorphic Cartesian square, upper submersion,
commuting projections and equal lengths. Whole actual marked-transform
comparison gives $N_Q=g^*M_T$, with no reindexing or filtering. Therefore
its entire boundary is the displayed literal list map and its final
cosupport is the inverse image of the empty target cosupport. The same
source history is thus terminal. Actual ambient canonical comparison
gives canonicality for the genuine source ambient dimension, including
an empty source. The source domain follows from actual smooth pullback.
The zero locus of its entire ideal is $b^{-1}F$, so the preceding lemma
supplies every stated source admissibility, principalization, support,
SNC, properness, surjectivity and isomorphism assertion. Two literal
list maps retain exactly the original number of positions. The target
was fixed before $b$, so no different target witness has been substituted.
\end{proof}

\begin{theorem}[An arbitrary compact set has a relatively compact canonical principalization neighbourhood]
\label{hir:original-local-canonical-principalization-compact-neighbourhood}
\lean{Hironaka.HironakaOriginalLocalCanonicalPrincipalization.exists_canonical_boundary_principalization_local}
\uses{hir:original-local-canonical-principalization-prescribed-open}
\leanok
For every original compact $K\subset X$ there is an original open
$U\supset K$ with compact closure and a genuine actual canonical
terminal tower for $M_0|_U$, retaining all target geometry and all
same-target arbitrary smooth-source conclusions of the preceding
theorem. The actual $U,T$ are chosen before all these source maps.
Neither a canonical choice of neighbourhood nor a comparison between
different neighbourhood choices is asserted.
\end{theorem}
\begin{proof}
The original complex manifold is locally compact. Compactness of $K$
therefore gives an original open neighbourhood $U$ of $K$ with compact
closure in $X$. Its whole set is contained in that closure. Apply the
preceding theorem with the prescribed $U$ and compact set $L=\overline U$.
Retain its actual target witness before all arbitrary smooth sources.
Every required target and source assertion is then exactly that theorem's
conclusion on this one neighbourhood. No compactness of the original
zero locus, source manifold or source cosupport enters this construction.
\end{proof}

For a global actual submersion $Y\to X$, its genuine restriction
$b^{-1}U\to U$ is an instance of the arbitrary source statement.
This local finite construction does not assert an unrestricted globally
locally finite resolution or a functorial rule for choosing $U$.
